\documentclass[a4paper,twoside]{article}

\usepackage{morewrites}

\usepackage[svgnames,dvipsnames,x11names]{xcolor}
\usepackage{graphicx}
\usepackage{texsmith-lists}
\usepackage[english]{babel}
\usepackage[colorlinks=true, linkcolor=NavyBlue, urlcolor=NavyBlue, citecolor=NavyBlue]{hyperref}
\usepackage[a4paper]{geometry}
\usepackage{ts-fonts}
\usepackage{amsmath}
\usepackage{float}
\usepackage{seqsplit}
\usepackage{ts-callouts}
\usepackage{ts-code}

\usepackage{lastpage}
\usepackage{refcount}
\makeatletter
\def\ps@plain{
\let\@oddhead\@empty
\let\@evenhead\@empty
\def\@oddfoot{
\hfil
\ifnum\getpagerefnumber{LastPage}>1\relax
\thepage
\fi
\hfil}
\let\@evenfoot\@oddfoot
}
\makeatother

\title{Fragments}
\author{}\date{}
\begin{document}
\maketitle
Fragments are reusable \LaTeX{} snippets injected into template slots. Built-in fragments (e.g., geometry, fonts, glossary, index) and custom ones share the same structure:
\begin{itemize}
\item \tscodeinline{fragment.toml} with either an \tscodeinline{entrypoint} --- the Python class then owns \tscodeinline{name}
and \tscodeinline{description} --- or a \tscodeinline{files} list, which declares them itself.
\item Optional \tscodeinline{attributes} section describing fragment-owned attributes (ownership enforced).
\item A \tscodeinline{provides} list naming the contract macros the fragment defines; a fragment that replaces another must provide every entry of its row.
\item Optional \tscodeinline{should\_render} logic (via entrypoint) to render only when needed.
\end{itemize}
\section{Manifest shape}
\begin{tscode}[lang=toml, engine=pygments]
\begin{Verbatim}[commandchars=\\\{\},breaklines, breakanywhere, commandchars=\\\{\}]
\PY{n}{name}\PY{+w}{ }\PY{o}{=}\PY{+w}{ }\PY{l+s+s2}{\PYZdq{}}\PY{l+s+s2}{my\PYZhy{}fragment}\PY{l+s+s2}{\PYZdq{}}
\PY{n}{description}\PY{+w}{ }\PY{o}{=}\PY{+w}{ }\PY{l+s+s2}{\PYZdq{}}\PY{l+s+s2}{Short description.}\PY{l+s+s2}{\PYZdq{}}

\PY{k}{[}\PY{k}{files}\PY{k}{.}\PY{k}{0}\PY{k}{]}
\PY{n}{path}\PY{+w}{ }\PY{o}{=}\PY{+w}{ }\PY{l+s+s2}{\PYZdq{}}\PY{l+s+s2}{my\PYZhy{}fragment.jinja.sty}\PY{l+s+s2}{\PYZdq{}}\PY{+w}{  }\PY{c+c1}{\PYZsh{} or .tex}
\PY{n}{slot}\PY{+w}{ }\PY{o}{=}\PY{+w}{ }\PY{l+s+s2}{\PYZdq{}}\PY{l+s+s2}{extra\PYZus{}packages}\PY{l+s+s2}{\PYZdq{}}
\PY{n}{type}\PY{+w}{ }\PY{o}{=}\PY{+w}{ }\PY{l+s+s2}{\PYZdq{}}\PY{l+s+s2}{package}\PY{l+s+s2}{\PYZdq{}}\PY{+w}{                }\PY{c+c1}{\PYZsh{} package | input | inline}
\PY{n}{output}\PY{+w}{ }\PY{o}{=}\PY{+w}{ }\PY{l+s+s2}{\PYZdq{}}\PY{l+s+s2}{my\PYZhy{}fragment}\PY{l+s+s2}{\PYZdq{}}

\PY{k}{[}\PY{k}{attributes}\PY{k}{.}\PY{k}{option}\PY{k}{]}
\PY{n}{default}\PY{+w}{ }\PY{o}{=}\PY{+w}{ }\PY{k+kc}{true}
\PY{n}{type}\PY{+w}{ }\PY{o}{=}\PY{+w}{ }\PY{l+s+s2}{\PYZdq{}}\PY{l+s+s2}{boolean}\PY{l+s+s2}{\PYZdq{}}
\PY{n}{sources}\PY{+w}{ }\PY{o}{=}\PY{+w}{ }\PY{p}{[}\PY{l+s+s2}{\PYZdq{}}\PY{l+s+s2}{press.option}\PY{l+s+s2}{\PYZdq{}}\PY{p}{,}\PY{+w}{ }\PY{l+s+s2}{\PYZdq{}}\PY{l+s+s2}{option}\PY{l+s+s2}{\PYZdq{}}\PY{p}{]}
\PY{n}{description}\PY{+w}{ }\PY{o}{=}\PY{+w}{ }\PY{l+s+s2}{\PYZdq{}}\PY{l+s+s2}{Controls feature X.}\PY{l+s+s2}{\PYZdq{}}
\end{Verbatim}
\end{tscode}
If you prefer Python logic, point \tscodeinline{entrypoint} at a callable returning a \tscodeinline{BaseFragment} instance.
\section{Fragment pieces}
\begin{itemize}
\item \tscodeinline{package}: rendered to a \tscodeinline{.sty} file and injected via \tscodeinline{\textbackslash{}usepackage} into the target slot.
\item \tscodeinline{input}: rendered to \tscodeinline{.tex} and injected via \tscodeinline{\textbackslash{}input\{\}} into the target slot.
\item \tscodeinline{inline}: rendered string injected directly into the slot variable (no file emitted).
\end{itemize}
Slots are validated against the template at render time; inline injections must match declared slots or template variables.
\section{Attributes}
Fragments can declare attributes (ownership enforced) and resolve overrides using the same \tscodeinline{TemplateAttributeSpec} model as templates. Attribute defaults are injected into the context before rendering fragment pieces. Attribute ownership matters: if two fragments (or a template) claim the same attribute name, TeXSmith raises a \tscodeinline{TemplateError}. Keep each attribute owned by a single fragment or the template to avoid conflicts.
\section{Contract macros}
A fragment defines the macros and environments the writers name, and says so:
\begin{tscode}[lang=toml, engine=pygments]
\begin{Verbatim}[commandchars=\\\{\},breaklines, breakanywhere, commandchars=\\\{\}]
\PY{n}{provides}\PY{+w}{ }\PY{o}{=}\PY{+w}{ }\PY{p}{[}\PY{l+s+s2}{\PYZdq{}}\PY{l+s+s2}{tscode}\PY{l+s+s2}{\PYZdq{}}\PY{p}{,}\PY{+w}{ }\PY{l+s+s2}{\PYZdq{}}\PY{l+s+s2}{tscodeinline}\PY{l+s+s2}{\PYZdq{}}\PY{p}{]}
\end{Verbatim}
\end{tscode}
The writer activates a fragment through \tscodeinline{Requires.fragments}: a body that emits
\tscodeinline{\textbackslash{}begin\{tscode\}} pulls in \tscodeinline{ts-\allowbreak{}code}, one that emits \tscodeinline{\textbackslash{}tsindex} pulls in
\tscodeinline{ts-\allowbreak{}index}. A macro the writer names, a fragment must define --- the check runs at
load time against \tscodeinline{tmark.fragments()}, so an incomplete replacement fails early
instead of producing an undefined control sequence.

Replace a bundled fragment outright with \tscodeinline{press.fragments}:
\begin{tscode}[lang=yaml, engine=pygments]
\begin{Verbatim}[commandchars=\\\{\},breaklines, breakanywhere, commandchars=\\\{\}]
\PY{n+nt}{press}\PY{p}{:}
\PY{+w}{  }\PY{n+nt}{fragments}\PY{p}{:}
\PY{+w}{    }\PY{n+nt}{disable}\PY{p}{:}\PY{+w}{ }\PY{p+pIndicator}{[}\PY{n+nv}{ts\PYZhy{}code}\PY{p+pIndicator}{]}
\PY{+w}{    }\PY{n+nt}{append}\PY{p}{:}\PY{+w}{ }\PY{p+pIndicator}{[}\PY{n+nv}{./my\PYZhy{}code.sty}\PY{p+pIndicator}{]}
\end{Verbatim}
\end{tscode}
See Contract macros for the table of macros and
their keys.
\begin{tscallout}[kind=warning, title={Jinja partials are no longer the rendering layer}]
A fragment's \tscodeinline{partials} and \tscodeinline{required\_partials} keys were deprecated in
0.7.0 and are removed in 0.8.0: \tscodeinline{fragment.toml} no longer reads them.
\end{tscallout}
\section{Runtime loading}
\begin{itemize}
\item Built-ins are discovered under \tscodeinline{src/texsmith/fragments/**/fragment.toml}.
\item Custom fragments can be passed via \tscodeinline{press.fragments} in front matter or CLI attributes.
\item The registry uses manifest metadata first; entrypoint is a fallback.
\item \textbf{Entrypoint contract (Python)}: Must expose a module attribute \tscodeinline{fragment} that is a \tscodeinline{BaseFragment} instance. Legacy \tscodeinline{create\_fragment()} factories are no longer used.
\end{itemize}
\section{Fragment contract (Python)}
\begin{itemize}
\item Implement a \tscodeinline{BaseFragment[Config]} subclass with:
\begin{itemize}
\item class attributes: \tscodeinline{name}, \tscodeinline{description}, \tscodeinline{pieces}, \tscodeinline{attributes} (TemplateAttributeSpec map), optional \tscodeinline{context\_defaults}, \tscodeinline{provides}, \tscodeinline{source}.
\item methods: \tscodeinline{build\_config(context, overrides=None) -\allowbreak{}> Config}, \tscodeinline{inject(config, context, overrides=None) -\allowbreak{}> None}, \tscodeinline{should\_render(config) -\allowbreak{}> bool}.
\end{itemize}
\item Define a \tscodeinline{Config} dataclass with \tscodeinline{from\_context(...)} and \tscodeinline{inject\_into(context)} methods (plus helpers like \tscodeinline{enabled()}).
\item Export \tscodeinline{fragment = YourFragment()} from \tscodeinline{\_\_init\_\_.py}; point \tscodeinline{fragment.toml} \tscodeinline{entrypoint} to \tscodeinline{texsmith.fragments.yourname:fragment}.
\item Injection flow: registry merges attribute defaults, builds config, calls \tscodeinline{inject}, then checks \tscodeinline{should\_render} before rendering pieces.
\end{itemize}
\subsection{Minimal example}
\tscodeinline{src/texsmith/fragments/example/\_\_init\_\_.py}
\begin{tscode}[lang=python, engine=pygments]
\begin{Verbatim}[commandchars=\\\{\},breaklines, breakanywhere, commandchars=\\\{\}]
\PY{k+kn}{from}\PY{+w}{ }\PY{n+nn}{dataclasses}\PY{+w}{ }\PY{k+kn}{import} \PY{n}{dataclass}
\PY{k+kn}{from}\PY{+w}{ }\PY{n+nn}{pathlib}\PY{+w}{ }\PY{k+kn}{import} \PY{n}{Path}
\PY{k+kn}{from}\PY{+w}{ }\PY{n+nn}{typing}\PY{+w}{ }\PY{k+kn}{import} \PY{n}{Any}\PY{p}{,} \PY{n}{Mapping}

\PY{k+kn}{from}\PY{+w}{ }\PY{n+nn}{texsmith}\PY{n+nn}{.}\PY{n+nn}{core}\PY{n+nn}{.}\PY{n+nn}{fragments}\PY{n+nn}{.}\PY{n+nn}{base}\PY{+w}{ }\PY{k+kn}{import} \PY{n}{BaseFragment}\PY{p}{,} \PY{n}{FragmentPiece}
\PY{k+kn}{from}\PY{+w}{ }\PY{n+nn}{texsmith}\PY{n+nn}{.}\PY{n+nn}{core}\PY{n+nn}{.}\PY{n+nn}{templates}\PY{n+nn}{.}\PY{n+nn}{manifest}\PY{+w}{ }\PY{k+kn}{import} \PY{n}{TemplateAttributeSpec}

\PY{n+nd}{@dataclass}\PY{p}{(}\PY{n}{frozen}\PY{o}{=}\PY{k+kc}{True}\PY{p}{)}
\PY{k}{class}\PY{+w}{ }\PY{n+nc}{ExampleConfig}\PY{p}{:}
    \PY{n}{message}\PY{p}{:} \PY{n+nb}{str} \PY{o}{|} \PY{k+kc}{None}

    \PY{n+nd}{@classmethod}
    \PY{k}{def}\PY{+w}{ }\PY{n+nf}{from\PYZus{}context}\PY{p}{(}\PY{n+nb+bp}{cls}\PY{p}{,} \PY{n}{ctx}\PY{p}{:} \PY{n}{Mapping}\PY{p}{[}\PY{n+nb}{str}\PY{p}{,} \PY{n}{Any}\PY{p}{]}\PY{p}{)} \PY{o}{\PYZhy{}}\PY{o}{\PYZgt{}} \PY{l+s+s2}{\PYZdq{}}\PY{l+s+s2}{ExampleConfig}\PY{l+s+s2}{\PYZdq{}}\PY{p}{:}
        \PY{k}{return} \PY{n+nb+bp}{cls}\PY{p}{(}\PY{n}{message}\PY{o}{=}\PY{n}{ctx}\PY{o}{.}\PY{n}{get}\PY{p}{(}\PY{l+s+s2}{\PYZdq{}}\PY{l+s+s2}{example\PYZus{}message}\PY{l+s+s2}{\PYZdq{}}\PY{p}{)}\PY{p}{)}

    \PY{k}{def}\PY{+w}{ }\PY{n+nf}{inject\PYZus{}into}\PY{p}{(}\PY{n+nb+bp}{self}\PY{p}{,} \PY{n}{ctx}\PY{p}{:} \PY{n+nb}{dict}\PY{p}{[}\PY{n+nb}{str}\PY{p}{,} \PY{n}{Any}\PY{p}{]}\PY{p}{)} \PY{o}{\PYZhy{}}\PY{o}{\PYZgt{}} \PY{k+kc}{None}\PY{p}{:}
        \PY{n}{ctx}\PY{p}{[}\PY{l+s+s2}{\PYZdq{}}\PY{l+s+s2}{ts\PYZus{}example\PYZus{}message}\PY{l+s+s2}{\PYZdq{}}\PY{p}{]} \PY{o}{=} \PY{n+nb+bp}{self}\PY{o}{.}\PY{n}{message} \PY{o+ow}{or} \PY{l+s+s2}{\PYZdq{}}\PY{l+s+s2}{Hello}\PY{l+s+s2}{\PYZdq{}}

    \PY{k}{def}\PY{+w}{ }\PY{n+nf}{enabled}\PY{p}{(}\PY{n+nb+bp}{self}\PY{p}{)} \PY{o}{\PYZhy{}}\PY{o}{\PYZgt{}} \PY{n+nb}{bool}\PY{p}{:}
        \PY{k}{return} \PY{n+nb}{bool}\PY{p}{(}\PY{n+nb+bp}{self}\PY{o}{.}\PY{n}{message}\PY{p}{)}

\PY{k}{class}\PY{+w}{ }\PY{n+nc}{ExampleFragment}\PY{p}{(}\PY{n}{BaseFragment}\PY{p}{[}\PY{n}{ExampleConfig}\PY{p}{]}\PY{p}{)}\PY{p}{:}
    \PY{n}{name} \PY{o}{=} \PY{l+s+s2}{\PYZdq{}}\PY{l+s+s2}{ts\PYZhy{}example}\PY{l+s+s2}{\PYZdq{}}
    \PY{n}{description} \PY{o}{=} \PY{l+s+s2}{\PYZdq{}}\PY{l+s+s2}{Tiny example fragment.}\PY{l+s+s2}{\PYZdq{}}
    \PY{n}{pieces} \PY{o}{=} \PY{p}{[}
        \PY{n}{FragmentPiece}\PY{p}{(}
            \PY{n}{template\PYZus{}path}\PY{o}{=}\PY{n}{Path}\PY{p}{(}\PY{n+nv+vm}{\PYZus{}\PYZus{}file\PYZus{}\PYZus{}}\PY{p}{)}\PY{o}{.}\PY{n}{with\PYZus{}name}\PY{p}{(}\PY{l+s+s2}{\PYZdq{}}\PY{l+s+s2}{ts\PYZhy{}example.jinja.tex}\PY{l+s+s2}{\PYZdq{}}\PY{p}{)}\PY{p}{,}
            \PY{n}{kind}\PY{o}{=}\PY{l+s+s2}{\PYZdq{}}\PY{l+s+s2}{inline}\PY{l+s+s2}{\PYZdq{}}\PY{p}{,}
            \PY{n}{slot}\PY{o}{=}\PY{l+s+s2}{\PYZdq{}}\PY{l+s+s2}{extra\PYZus{}packages}\PY{l+s+s2}{\PYZdq{}}\PY{p}{,}
        \PY{p}{)}
    \PY{p}{]}
    \PY{n}{attributes} \PY{o}{=} \PY{p}{\PYZob{}}
        \PY{l+s+s2}{\PYZdq{}}\PY{l+s+s2}{example\PYZus{}message}\PY{l+s+s2}{\PYZdq{}}\PY{p}{:} \PY{n}{TemplateAttributeSpec}\PY{p}{(}
            \PY{n}{default}\PY{o}{=}\PY{k+kc}{None}\PY{p}{,}
            \PY{n+nb}{type}\PY{o}{=}\PY{l+s+s2}{\PYZdq{}}\PY{l+s+s2}{string}\PY{l+s+s2}{\PYZdq{}}\PY{p}{,}
            \PY{n}{sources}\PY{o}{=}\PY{p}{[}\PY{l+s+s2}{\PYZdq{}}\PY{l+s+s2}{press.example.message}\PY{l+s+s2}{\PYZdq{}}\PY{p}{,} \PY{l+s+s2}{\PYZdq{}}\PY{l+s+s2}{example.message}\PY{l+s+s2}{\PYZdq{}}\PY{p}{]}\PY{p}{,}
        \PY{p}{)}
    \PY{p}{\PYZcb{}}
    \PY{n}{context\PYZus{}defaults}\PY{p}{:} \PY{n+nb}{dict}\PY{p}{[}\PY{n+nb}{str}\PY{p}{,} \PY{n}{Any}\PY{p}{]} \PY{o}{=} \PY{p}{\PYZob{}}\PY{l+s+s2}{\PYZdq{}}\PY{l+s+s2}{extra\PYZus{}packages}\PY{l+s+s2}{\PYZdq{}}\PY{p}{:} \PY{l+s+s2}{\PYZdq{}}\PY{l+s+s2}{\PYZdq{}}\PY{p}{\PYZcb{}}
    \PY{n}{config\PYZus{}cls} \PY{o}{=} \PY{n}{ExampleConfig}
    \PY{n}{source} \PY{o}{=} \PY{n}{Path}\PY{p}{(}\PY{n+nv+vm}{\PYZus{}\PYZus{}file\PYZus{}\PYZus{}}\PY{p}{)}\PY{o}{.}\PY{n}{with\PYZus{}name}\PY{p}{(}\PY{l+s+s2}{\PYZdq{}}\PY{l+s+s2}{ts\PYZhy{}example.jinja.tex}\PY{l+s+s2}{\PYZdq{}}\PY{p}{)}

    \PY{k}{def}\PY{+w}{ }\PY{n+nf}{build\PYZus{}config}\PY{p}{(}\PY{n+nb+bp}{self}\PY{p}{,} \PY{n}{context}\PY{p}{:} \PY{n}{Mapping}\PY{p}{[}\PY{n+nb}{str}\PY{p}{,} \PY{n}{Any}\PY{p}{]}\PY{p}{,} \PY{n}{overrides}\PY{o}{=}\PY{k+kc}{None}\PY{p}{)} \PY{o}{\PYZhy{}}\PY{o}{\PYZgt{}} \PY{n}{ExampleConfig}\PY{p}{:}
        \PY{k}{return} \PY{n+nb+bp}{self}\PY{o}{.}\PY{n}{config\PYZus{}cls}\PY{o}{.}\PY{n}{from\PYZus{}context}\PY{p}{(}\PY{n}{context}\PY{p}{)}

    \PY{k}{def}\PY{+w}{ }\PY{n+nf}{inject}\PY{p}{(}\PY{n+nb+bp}{self}\PY{p}{,} \PY{n}{config}\PY{p}{:} \PY{n}{ExampleConfig}\PY{p}{,} \PY{n}{context}\PY{p}{:} \PY{n+nb}{dict}\PY{p}{[}\PY{n+nb}{str}\PY{p}{,} \PY{n}{Any}\PY{p}{]}\PY{p}{,} \PY{n}{overrides}\PY{o}{=}\PY{k+kc}{None}\PY{p}{)} \PY{o}{\PYZhy{}}\PY{o}{\PYZgt{}} \PY{k+kc}{None}\PY{p}{:}
        \PY{n}{config}\PY{o}{.}\PY{n}{inject\PYZus{}into}\PY{p}{(}\PY{n}{context}\PY{p}{)}

    \PY{k}{def}\PY{+w}{ }\PY{n+nf}{should\PYZus{}render}\PY{p}{(}\PY{n+nb+bp}{self}\PY{p}{,} \PY{n}{config}\PY{p}{:} \PY{n}{ExampleConfig}\PY{p}{)} \PY{o}{\PYZhy{}}\PY{o}{\PYZgt{}} \PY{n+nb}{bool}\PY{p}{:}
        \PY{k}{return} \PY{n}{config}\PY{o}{.}\PY{n}{enabled}\PY{p}{(}\PY{p}{)}

\PY{n}{fragment} \PY{o}{=} \PY{n}{ExampleFragment}\PY{p}{(}\PY{p}{)}
\end{Verbatim}
\end{tscode}
\tscodeinline{src/texsmith/fragments/example/fragment.toml}
\begin{tscode}[lang=toml, engine=pygments]
\begin{Verbatim}[commandchars=\\\{\},breaklines, breakanywhere, commandchars=\\\{\}]
\PY{n}{name}\PY{+w}{ }\PY{o}{=}\PY{+w}{ }\PY{l+s+s2}{\PYZdq{}}\PY{l+s+s2}{ts\PYZhy{}example}\PY{l+s+s2}{\PYZdq{}}
\PY{n}{description}\PY{+w}{ }\PY{o}{=}\PY{+w}{ }\PY{l+s+s2}{\PYZdq{}}\PY{l+s+s2}{Tiny example fragment.}\PY{l+s+s2}{\PYZdq{}}
\PY{n}{entrypoint}\PY{+w}{ }\PY{o}{=}\PY{+w}{ }\PY{l+s+s2}{\PYZdq{}}\PY{l+s+s2}{texsmith.fragments.example:fragment}\PY{l+s+s2}{\PYZdq{}}

\PY{k}{[[}\PY{k}{files}\PY{k}{]]}
\PY{n}{path}\PY{+w}{ }\PY{o}{=}\PY{+w}{ }\PY{l+s+s2}{\PYZdq{}}\PY{l+s+s2}{ts\PYZhy{}example.jinja.tex}\PY{l+s+s2}{\PYZdq{}}
\PY{n}{slot}\PY{+w}{ }\PY{o}{=}\PY{+w}{ }\PY{l+s+s2}{\PYZdq{}}\PY{l+s+s2}{extra\PYZus{}packages}\PY{l+s+s2}{\PYZdq{}}
\PY{n}{type}\PY{+w}{ }\PY{o}{=}\PY{+w}{ }\PY{l+s+s2}{\PYZdq{}}\PY{l+s+s2}{inline}\PY{l+s+s2}{\PYZdq{}}
\end{Verbatim}
\end{tscode}
\section{Migration notes for fragment authors}
\begin{itemize}
\item Prefer \tscodeinline{fragment.toml} with \tscodeinline{files}, \tscodeinline{attributes}, and \tscodeinline{provides}; use an entrypoint only when you need Python logic.
\item If using Python, return a \tscodeinline{BaseFragment} instance via \tscodeinline{fragment}; \tscodeinline{create\_fragment()} and \tscodeinline{FragmentDefinition} shims have been removed.
\item Declare attribute ownership (implicit \tscodeinline{owner = <fragment name>}) to avoid collisions with templates or other fragments.
\item Keep slot targets aligned with template variables; inline targets must reference declared slots or template variables, otherwise a \tscodeinline{TemplateError} is raised.
\end{itemize}
\end{document}
